\clearpage
\section{교대군과 대칭군의 정상구조}
\label{sec:alternating-symmetric-structure}

\begin{learninggoal}
$n\ge5$에서 $A_n$이 단순군임을 완전하게 증명하고, 그 결과로 $S_n$의 정상부분군과 도함열을 분류한다. 비가환 단순군이 가해군의 정규열에 들어갈 수 없는 이유를 명확히 한다.
\end{learninggoal}

$A_5$의 켤레류 계산은 오차방정식에 필요한 가장 구체적인 단순성 증명이다. 그러나 같은 현상은 모든 $A_n$, $n\ge5$에서 성립한다. 이 절에서는 삼순환과 교환자를 사용하여 이를 증명한다.

\subsection{삼순환이 만드는 교대군}

\begin{proposition}
\label{prop:an-generated-three-cycles}
$n\ge3$이면 $A_n$은 삼순환들로 생성된다.
\end{proposition}

\begin{proof}
모든 짝순열은 짝수 개의 전치의 곱이다. 전치 두 개는 다음 두 유형으로 처리할 수 있다.
\[
  (a\,b)(b\,c)=(a\,b\,c),
\]
서로소인 경우에는
\[
  (a\,b)(c\,d)=(a\,c\,b)(a\,c\,d).
\]
따라서 전치들을 두 개씩 묶으면 모든 짝순열을 삼순환들의 곱으로 쓸 수 있다.
\end{proof}

\begin{lemma}
\label{lem:three-cycles-conjugate-an}
$n\ge5$이면 $A_n$의 모든 삼순환은 $A_n$ 안에서 서로 켤레이다.
\end{lemma}

\begin{proof}
$S_n$에서는 같은 순환형의 원소들이 서로 켤레이므로 두 삼순환 $\alpha,\beta$에 대해
\[
  \pi\alpha\pi^{-1}=\beta
\]
인 $\pi\in S_n$이 존재한다. $\pi$가 짝순열이면 끝이다. $\pi$가 홀순열이면 $\beta$가 움직이는 세 점 밖에 서로 다른 두 점 $u,v$를 택할 수 있다. 전치 $(u\,v)$는 $\beta$와 가환하므로
\[
  ((u\,v)\pi)\alpha((u\,v)\pi)^{-1}=\beta,
\]
그리고 $(u\,v)\pi$는 짝순열이다.
\end{proof}

\subsection{교대군의 단순성}

\begin{theorem}[교대군의 단순성]
\label{thm:an-simple}
$n\ge5$이면 $A_n$은 단순군이다.
\end{theorem}

\begin{proofidea}
비자명한 정상부분군 $N$에서 원소 $\sigma\ne1$을 하나 잡는다. 적절한 삼순환 $\tau$와의 교환자
\[
  [\tau,\sigma]=\tau\sigma\tau^{-1}\sigma^{-1}
\]
를 계산하여 $N$ 안에 삼순환을 만든다. 정상성 때문에 $N$은 모든 삼순환을 포함하고, 삼순환들이 $A_n$을 생성하므로 $N=A_n$이다.
\end{proofidea}

\begin{proof}
$1\ne N\triangleleft A_n$이라 하고 $1\ne\sigma\in N$을 택하자. $N$은 정상부분군이므로 모든 $\tau\in A_n$에 대해
\[
  [\tau,\sigma]=\tau\sigma\tau^{-1}\sigma^{-1}\in N
\]
이다. $\sigma$의 서로소 순환분해에 따라 경우를 나눈다.

\medskip
\noindent\textbf{경우 1: $\sigma$가 길이 $4$ 이상인 순환을 포함한다.}
그 순환의 연속한 네 점을 $a,b,c,d$라 쓰고
\[
  \tau=(a\,b\,c)
\]
라 하자. 직접 계산하면
\[
  [\tau,\sigma]=(a\,b\,d)
\]
이다. 따라서 $N$은 삼순환을 포함한다.

\medskip
\noindent\textbf{경우 2: $\sigma$가 삼순환 $(a\,b\,c)$을 포함한다.}
$\sigma=(a\,b\,c)$이면 이미 끝이다. 그렇지 않으면 $\{a,b,c\}$ 밖에서 $\sigma$가 움직이는 점 $d$를 택하고 $e=\sigma(d)$라 하자. 서로소 순환분해 때문에 $d,e$는 $a,b,c$와 다르다. 이제
\[
  \tau=(a\,b\,d)
\]
라 두면 직접 계산으로
\[
  [\tau,\sigma]=(a\,b\,e\,c\,d)
\]
인 오순환을 얻는다. 이 원소는 $N$에 속하고 길이가 $4$ 이상인 순환이므로 경우 1을 다시 적용하면 $N$ 안에 삼순환이 존재한다.

\medskip
\noindent\textbf{경우 3: $\sigma$는 서로소인 전치들의 곱이다.}
$\sigma\in A_n$이므로 전치의 개수는 짝수이고 적어도 두 개이다.

먼저 $\sigma$가 어떤 점 $e$를 고정한다고 하자. 전치 하나를 $(a\,b)$라 쓰고 $\tau=(a\,b\,e)$로 두면
\[
  [\tau,\sigma]=(a\,e\,b)
\]
이므로 $N$은 삼순환을 포함한다.

이제 $\sigma$가 고정점을 갖지 않는다고 하자. $n\ge5$이고 전치의 개수가 짝수이므로 적어도 네 개의 서로소 전치가 있다. 그중 세 개를
\[
  (a\,b)(c\,d)(e\,f)
\]
라 쓰고 $\tau=(a\,c\,e)$로 두면
\[
  [\tau,\sigma]=(a\,c\,e)(b\,f\,d)
\]
이다. 이 원소는 두 삼순환의 곱이므로 경우 2를 적용하여 다시 $N$ 안에 삼순환을 얻는다.

세 경우 모두 $N$은 어떤 삼순환을 포함한다. 보조정리 \ref{lem:three-cycles-conjugate-an}와 정상성에 의해 $N$은 모든 삼순환을 포함한다. 명제 \ref{prop:an-generated-three-cycles}에서 삼순환들이 $A_n$을 생성하므로
\[
  N=A_n.
\]
\end{proof}

\begin{remark}
$n=4$에서는 정리가 거짓이다. 실제로
\[
  V_4\triangleleft A_4
\]
는 비자명한 진정한 정상부분군이다. 따라서 $n=5$는 교대군의 단순성이 시작되는 정확한 경계이다.
\end{remark}

\subsection{대칭군의 정상부분군과 도함열}

\begin{theorem}
\label{thm:sn-normal-subgroups}
$n\ge5$일 때 $S_n$의 정상부분군은
\[
  1,\qquad A_n,\qquad S_n
\]
뿐이다.
\end{theorem}

\begin{proof}
$N\triangleleft S_n$이라 하자. 그러면
\[
  N\cap A_n\triangleleft A_n.
\]
정리 \ref{thm:an-simple}에 의해 $N\cap A_n$은 $1$ 또는 $A_n$이다.

$N\cap A_n=A_n$이면 $A_n\subseteq N$이고 $[S_n:A_n]=2$이므로 $N=A_n$ 또는 $S_n$이다.

이제 $N\cap A_n=1$이라 하자. 부호준동형의 $N$에 대한 제한
\[
  \sgn|_N:N\longrightarrow\{\pm1\}
\]
은 단사이므로 $|N|\le2$이다. $N$이 비자명하면 $N=\{1,\sigma\}$이고, 정상성 때문에 모든 $g\in S_n$에 대해 $g\sigma g^{-1}=\sigma$이다. 즉 $\sigma$는 $S_n$의 중심에 속한다. 그러나 $n\ge3$에서 $S_n$의 중심은 자명하므로 모순이다. 따라서 $N=1$이다.
\end{proof}

\begin{definition}[완전군]
군 $G$가
\[
  [G,G]=G
\]
를 만족하면 $G$를 \emph{완전군}(perfect group)이라 한다.
\end{definition}

\begin{proposition}
\label{prop:sn-an-derived-series}
$n\ge5$이면
\[
  [S_n,S_n]=A_n,
  \qquad
  [A_n,A_n]=A_n.
\]
따라서
\[
  S_n^{(0)}=S_n,
  \qquad
  S_n^{(r)}=A_n\quad(r\ge1),
\]
이고 $S_n$과 $A_n$은 가해군이 아니다.
\end{proposition}

\begin{proof}
부호준동형의 핵이 $A_n$이므로 $S_n/A_n$은 아벨군이고
\[
  [S_n,S_n]\subseteq A_n.
\]
반대로 모든 삼순환은 두 전치의 교환자로 쓸 수 있고 삼순환들이 $A_n$을 생성하므로 $A_n\subseteq[S_n,S_n]$이다.

또한 $n\ge5$에서 삼순환 $(a\,b\,c)$에 대해 서로 다른 $d,e$를 택하면
\[
  (a\,b\,c)
  =[(a\,b\,d),(a\,c\,e)]
\]
와 같은 교환자 표현을 얻는다. 두 원소는 모두 $A_n$에 속하므로 모든 삼순환이 $[A_n,A_n]$에 들어간다. 따라서
\[
  [A_n,A_n]=A_n.
\]
도함열이 $A_n$에서 멈추므로 두 군은 가해군이 아니다.
\end{proof}

\begin{pitfall}
단순군과 완전군은 다른 개념이다. 소수차 순환군 $C_p$는 단순하지만 아벨군이므로 완전군이 아니다. 반면 $A_5$는 비가환 단순군이고 완전군이다. 근호해법을 막는 직접적인 조건은 도함열이 끝나지 않는 완전성이고, 단순성은 그 완전한 핵을 더 이상 정상층으로 나눌 수 없음을 보여 준다.
\end{pitfall}

\begin{checkpoint}
\begin{enumerate}[label=(\alph*)]
  \item 정리 \ref{thm:an-simple}의 경우 1에서 교환자를 직접 계산하라.
  \item $A_4$의 도함열을 적고, 왜 $A_4$는 단순하지 않지만 가해군인지 설명하라.
  \item $n\ge5$에서 $S_n/A_n\cong C_2$라는 사실만으로 $S_n$의 비가해성을 결론 낼 수 없는 이유를 설명하라.
\end{enumerate}
\end{checkpoint}
